\documentclass[letterpaper]{article} % DO NOT CHANGE THIS
\usepackage[submission]{aaai2027}  % DO NOT CHANGE THIS
\usepackage[hyphens]{url}  % DO NOT CHANGE THIS
\usepackage{graphicx} % DO NOT CHANGE THIS
\urlstyle{rm} % DO NOT CHANGE THIS
\def\UrlFont{\rm}  % DO NOT CHANGE THIS
\usepackage{natbib}  % DO NOT CHANGE THIS AND DO NOT ADD ANY OPTIONS TO IT
\usepackage{caption} % DO NOT CHANGE THIS AND DO NOT ADD ANY OPTIONS TO IT
\frenchspacing  % DO NOT CHANGE THIS
\usepackage{algorithm}
\usepackage{algorithmic}
\usepackage{booktabs}
\usepackage{amsmath}
\usepackage{amsthm}
\pdfinfo{
/TemplateVersion (2027.1)
}
% Every result below is a macro generated by make_numbers.py (claims.json);
% check_claims.py verifies that no result number is typed by hand.
% generated by make_numbers.py -- do not edit
\newcommand{\AuditInvalid}{0}
\newcommand{\AuditSolutions}{977}
\newcommand{\DiagDecisive}{3}
\newcommand{\DiagExplained}{3}
\newcommand{\DiagInfeasible}{29}
\newcommand{\DiagInfeasibleMs}{21}
\newcommand{\DiagOptimalS}{0.52}
\newcommand{\DiagSteps}{120}
\newcommand{\OracleCpsatCases}{1{,}000}
\newcommand{\OracleCpsatDecisive}{358}
\newcommand{\OracleCpsatViolations}{0}
\newcommand{\OracleCutCases}{1{,}000}
\newcommand{\OracleCutMisses}{0}
\newcommand{\OracleEnumCases}{2{,}000}
\newcommand{\OracleEnumDecisive}{808}
\newcommand{\OracleEnumViolations}{0}
\newcommand{\OracleGlobalCases}{500}
\newcommand{\OracleGlobalDecisive}{82}
\newcommand{\OracleGlobalViolations}{0}
\newcommand{\PilotCertDecisive}{13}
\newcommand{\PilotCertMs}{1.3}
\newcommand{\PilotCertSteps}{200}
\newcommand{\PilotEightASixrVsAOneCI}{[-0.0060, +0.0005]}
\newcommand{\PilotEightASixrVsAOneHolm}{1.00}
\newcommand{\PilotEightASixrVsAOneRel}{11.1}
\newcommand{\PilotEightASixrVsAOneSign}{-}
\newcommand{\PilotEightASixrVsAOneWL}{9/11}
\newcommand{\PilotEightCOneVsAOneCI}{[-0.0059, +0.0011]}
\newcommand{\PilotEightCOneVsAOneHolm}{1.00}
\newcommand{\PilotEightCOneVsAOneRel}{9.9}
\newcommand{\PilotEightCOneVsAOneSign}{-}
\newcommand{\PilotEightCOneVsAOneWL}{10/10}
\newcommand{\PilotEightCOneVsCThreeCI}{[-0.0059, +0.0017]}
\newcommand{\PilotEightCOneVsCThreeHolm}{1.00}
\newcommand{\PilotEightCOneVsCThreeRel}{8.8}
\newcommand{\PilotEightCOneVsCThreeSign}{-}
\newcommand{\PilotEightCOneVsCThreeWL}{11/9}
\newcommand{\PilotEightCThreeVsAOneCI}{[-0.0031, +0.0024]}
\newcommand{\PilotEightCThreeVsAOneHolm}{1.00}
\newcommand{\PilotEightCThreeVsAOneRel}{1.2}
\newcommand{\PilotEightCThreeVsAOneSign}{-}
\newcommand{\PilotEightCThreeVsAOneWL}{10/10}
\newcommand{\PilotEightCTwoVsAOneCI}{[-0.0124, -0.0036]}
\newcommand{\PilotEightCTwoVsAOneHolm}{0.095}
\newcommand{\PilotEightCTwoVsAOneRel}{33.5}
\newcommand{\PilotEightCTwoVsAOneSign}{-}
\newcommand{\PilotEightCTwoVsAOneWL}{16/4}
\newcommand{\PilotEightCTwoVsCThreeCI}{[-0.0121, -0.0036]}
\newcommand{\PilotEightCTwoVsCThreeHolm}{0.095}
\newcommand{\PilotEightCTwoVsCThreeRel}{32.7}
\newcommand{\PilotEightCTwoVsCThreeSign}{-}
\newcommand{\PilotEightCTwoVsCThreeWL}{16/4}
\newcommand{\PilotFourAFourVsAOneCI}{[-0.0005, +0.0056]}
\newcommand{\PilotFourAFourVsAOneHolm}{0.55}
\newcommand{\PilotFourAFourVsAOneRel}{23.7}
\newcommand{\PilotFourAFourVsAOneSign}{+}
\newcommand{\PilotFourAFourVsAOneWL}{2/8}
\newcommand{\PilotFourAOneVsAZeroCI}{[-0.0069, -0.0020]}
\newcommand{\PilotFourAOneVsAZeroHolm}{0.13}
\newcommand{\PilotFourAOneVsAZeroRel}{27.5}
\newcommand{\PilotFourAOneVsAZeroSign}{-}
\newcommand{\PilotFourAOneVsAZeroWL}{9/1}
\newcommand{\PilotLocalMeanS}{0.48}
\newcommand{\PilotOneAFiveVsAFourCI}{[-0.0010, +0.0021]}
\newcommand{\PilotOneAFiveVsAFourHolm}{1.00}
\newcommand{\PilotOneAFiveVsAFourRel}{2.2}
\newcommand{\PilotOneAFiveVsAFourSign}{+}
\newcommand{\PilotOneAFiveVsAFourWL}{11/9}
\newcommand{\PilotOneAFourVsAOneCI}{[-0.0023, +0.0000]}
\newcommand{\PilotOneAFourVsAOneHolm}{1.00}
\newcommand{\PilotOneAFourVsAOneRel}{4.1}
\newcommand{\PilotOneAFourVsAOneSign}{-}
\newcommand{\PilotOneAFourVsAOneWL}{11/9}
\newcommand{\PilotOneANineVsAZeroCI}{[-0.0033, +0.0023]}
\newcommand{\PilotOneANineVsAZeroHolm}{1.00}
\newcommand{\PilotOneANineVsAZeroRel}{1.4}
\newcommand{\PilotOneANineVsAZeroSign}{-}
\newcommand{\PilotOneANineVsAZeroWL}{9/11}
\newcommand{\PilotOneASixVsAFourCI}{[-0.0011, +0.0022]}
\newcommand{\PilotOneASixVsAFourHolm}{1.00}
\newcommand{\PilotOneASixVsAFourRel}{2.4}
\newcommand{\PilotOneASixVsAFourSign}{+}
\newcommand{\PilotOneASixVsAFourWL}{12/8}
\newcommand{\PilotOneASixrVsAOneCI}{[-0.0095, -0.0029]}
\newcommand{\PilotOneASixrVsAOneHolm}{1.00}
\newcommand{\PilotOneASixrVsAOneRel}{26.2}
\newcommand{\PilotOneASixrVsAOneSign}{-}
\newcommand{\PilotOneASixrVsAOneWL}{14/6}
\newcommand{\PilotOneAThreeVsAOneCI}{[-0.0009, +0.0008]}
\newcommand{\PilotOneAThreeVsAOneHolm}{1.00}
\newcommand{\PilotOneAThreeVsAOneRel}{0.4}
\newcommand{\PilotOneAThreeVsAOneSign}{-}
\newcommand{\PilotOneAThreeVsAOneWL}{10/10}
\newcommand{\PilotOneATwoVsAOneCI}{[-0.0013, +0.0003]}
\newcommand{\PilotOneATwoVsAOneHolm}{1.00}
\newcommand{\PilotOneATwoVsAOneRel}{1.4}
\newcommand{\PilotOneATwoVsAOneSign}{-}
\newcommand{\PilotOneATwoVsAOneWL}{8/12}
\newcommand{\PilotReleaseMs}{1.5}
\newcommand{\PilotReleaseNeeded}{17}
\newcommand{\PilotShortStage}{62}
\newcommand{\PilotTenCTwoVsAOneCI}{[-0.0101, -0.0029]}
\newcommand{\PilotTenCTwoVsAOneHolm}{0.047}
\newcommand{\PilotTenCTwoVsAOneRel}{28.3}
\newcommand{\PilotTenCTwoVsAOneSign}{-}
\newcommand{\PilotTenCTwoVsAOneWL}{16/4}
\newcommand{\PilotTenCTwoVsCThreeCI}{[-0.0103, -0.0021]}
\newcommand{\PilotTenCTwoVsCThreeHolm}{0.12}
\newcommand{\PilotTenCTwoVsCThreeRel}{26.7}
\newcommand{\PilotTenCTwoVsCThreeSign}{-}
\newcommand{\PilotTenCTwoVsCThreeWL}{15/5}
\newcommand{\PilotTwoAFourVsAOneCI}{[+0.0023, +0.0181]}
\newcommand{\PilotTwoAFourVsAOneHolm}{1.00}
\newcommand{\PilotTwoAFourVsAOneRel}{41.0}
\newcommand{\PilotTwoAFourVsAOneSign}{+}
\newcommand{\PilotTwoAFourVsAOneWL}{3/7}
\newcommand{\PilotTwoAOneVsAZeroCI}{[-0.0185, -0.0040]}
\newcommand{\PilotTwoAOneVsAZeroHolm}{0.88}
\newcommand{\PilotTwoAOneVsAZeroRel}{31.3}
\newcommand{\PilotTwoAOneVsAZeroSign}{-}
\newcommand{\PilotTwoAOneVsAZeroWL}{8/2}
\newcommand{\SFiveCandidateVsIaipSrCI}{[-0.0054, -0.0015]}
\newcommand{\SFiveCandidateVsIaipSrGapCI}{[-0.0035, -0.0003]}
\newcommand{\SFiveCandidateVsIaipSrGapRel}{16.3}
\newcommand{\SFiveCandidateVsIaipSrGapSign}{-}
\newcommand{\SFiveCandidateVsIaipSrGapWL}{15/4}
\newcommand{\SFiveCandidateVsIaipSrHolm}{0.012}
\newcommand{\SFiveCandidateVsIaipSrRel}{14.7}
\newcommand{\SFiveCandidateVsIaipSrSign}{-}
\newcommand{\SFiveCandidateVsIaipSrWL}{16/4}
\newcommand{\SFiveCandidateVsPortfolioHintedCI}{[-0.0103, -0.0051]}
\newcommand{\SFiveCandidateVsPortfolioHintedGapCI}{[-0.0052, -0.0020]}
\newcommand{\SFiveCandidateVsPortfolioHintedGapRel}{26.3}
\newcommand{\SFiveCandidateVsPortfolioHintedGapSign}{-}
\newcommand{\SFiveCandidateVsPortfolioHintedGapWL}{17/2}
\newcommand{\SFiveCandidateVsPortfolioHintedHolm}{3.8\cdot10^{-6}}
\newcommand{\SFiveCandidateVsPortfolioHintedRel}{27.7}
\newcommand{\SFiveCandidateVsPortfolioHintedSign}{-}
\newcommand{\SFiveCandidateVsPortfolioHintedWL}{20/0}
\newcommand{\SFiveGapCandidate}{.0099}
\newcommand{\SFiveGapIaipSr}{.0118}
\newcommand{\SFiveGapPortfolioHinted}{.0134}
\newcommand{\SFiveMeanCandidate}{.0199}
\newcommand{\SFiveMeanIaipSr}{.0233}
\newcommand{\SFiveMeanPortfolioHinted}{.0275}
\newcommand{\SFiveNonInferior}{yes}
\newcommand{\SFiveNotValidated}{0}
\newcommand{\SFiveRecords}{3{,}300}
\newcommand{\SFiveUnits}{20}
\newcommand{\SFourCandidateVsIaipSrCI}{[-0.0225, -0.0130]}
\newcommand{\SFourCandidateVsIaipSrGapCI}{[-0.0227, -0.0127]}
\newcommand{\SFourCandidateVsIaipSrGapRel}{56.3}
\newcommand{\SFourCandidateVsIaipSrGapSign}{-}
\newcommand{\SFourCandidateVsIaipSrGapWL}{20/0}
\newcommand{\SFourCandidateVsIaipSrHolm}{3.8\cdot10^{-6}}
\newcommand{\SFourCandidateVsIaipSrRel}{40.8}
\newcommand{\SFourCandidateVsIaipSrSign}{-}
\newcommand{\SFourCandidateVsIaipSrWL}{20/0}
\newcommand{\SFourCandidateVsPortfolioHintedCI}{[-0.0523, -0.0410]}
\newcommand{\SFourCandidateVsPortfolioHintedGapCI}{[-0.0459, -0.0324]}
\newcommand{\SFourCandidateVsPortfolioHintedGapRel}{74.3}
\newcommand{\SFourCandidateVsPortfolioHintedGapSign}{-}
\newcommand{\SFourCandidateVsPortfolioHintedGapWL}{20/0}
\newcommand{\SFourCandidateVsPortfolioHintedHolm}{3.8\cdot10^{-6}}
\newcommand{\SFourCandidateVsPortfolioHintedRel}{64.8}
\newcommand{\SFourCandidateVsPortfolioHintedSign}{-}
\newcommand{\SFourCandidateVsPortfolioHintedWL}{20/0}
\newcommand{\SFourGapCandidate}{.0135}
\newcommand{\SFourGapIaipSr}{.0309}
\newcommand{\SFourGapPortfolioHinted}{.0525}
\newcommand{\SFourMeanCandidate}{.0253}
\newcommand{\SFourMeanIaipSr}{.0428}
\newcommand{\SFourMeanPortfolioHinted}{.0719}
\newcommand{\SFourNonInferior}{yes}
\newcommand{\SFourNotValidated}{0}
\newcommand{\SFourRecords}{3{,}300}
\newcommand{\SFourUnits}{20}
\newcommand{\SOneCandidateVsIaipSrCI}{[-0.0108, -0.0063]}
\newcommand{\SOneCandidateVsIaipSrGapCI}{[-0.0096, -0.0050]}
\newcommand{\SOneCandidateVsIaipSrGapRel}{31.1}
\newcommand{\SOneCandidateVsIaipSrGapSign}{-}
\newcommand{\SOneCandidateVsIaipSrGapWL}{37/3}
\newcommand{\SOneCandidateVsIaipSrHolm}{2.2\cdot10^{-10}}
\newcommand{\SOneCandidateVsIaipSrRel}{24.8}
\newcommand{\SOneCandidateVsIaipSrSign}{-}
\newcommand{\SOneCandidateVsIaipSrWL}{39/1}
\newcommand{\SOneCandidateVsPortfolioHintedCI}{[-0.0222, -0.0161]}
\newcommand{\SOneCandidateVsPortfolioHintedGapCI}{[-0.0159, -0.0099]}
\newcommand{\SOneCandidateVsPortfolioHintedGapRel}{44.5}
\newcommand{\SOneCandidateVsPortfolioHintedGapSign}{-}
\newcommand{\SOneCandidateVsPortfolioHintedGapWL}{38/2}
\newcommand{\SOneCandidateVsPortfolioHintedHolm}{7.3\cdot10^{-12}}
\newcommand{\SOneCandidateVsPortfolioHintedRel}{42.7}
\newcommand{\SOneCandidateVsPortfolioHintedSign}{-}
\newcommand{\SOneCandidateVsPortfolioHintedWL}{40/0}
\newcommand{\SOneCandidateVsRandomRoundsCI}{[-0.0073, -0.0039]}
\newcommand{\SOneCandidateVsRandomRoundsGapCI}{[-0.0066, -0.0029]}
\newcommand{\SOneCandidateVsRandomRoundsGapRel}{22.7}
\newcommand{\SOneCandidateVsRandomRoundsGapSign}{-}
\newcommand{\SOneCandidateVsRandomRoundsGapWL}{34/6}
\newcommand{\SOneCandidateVsRandomRoundsHolm}{3.0\cdot10^{-9}}
\newcommand{\SOneCandidateVsRandomRoundsRel}{17.7}
\newcommand{\SOneCandidateVsRandomRoundsSign}{-}
\newcommand{\SOneCandidateVsRandomRoundsWL}{38/2}
\newcommand{\SOneDeltaArrivalVsIaipSrCI}{[-0.0116, -0.0071]}
\newcommand{\SOneDeltaArrivalVsIaipSrRel}{25.4}
\newcommand{\SOneDeltaArrivalVsIaipSrSign}{-}
\newcommand{\SOneDeltaArrivalVsIaipSrWL}{38/2}
\newcommand{\SOneDeltaCancellationVsIaipSrCI}{[-0.0093, -0.0033]}
\newcommand{\SOneDeltaCancellationVsIaipSrRel}{19.7}
\newcommand{\SOneDeltaCancellationVsIaipSrSign}{-}
\newcommand{\SOneDeltaCancellationVsIaipSrWL}{28/7}
\newcommand{\SOneDeltaDurationJitterVsIaipSrCI}{[-0.0110, -0.0053]}
\newcommand{\SOneDeltaDurationJitterVsIaipSrRel}{25.8}
\newcommand{\SOneDeltaDurationJitterVsIaipSrSign}{-}
\newcommand{\SOneDeltaDurationJitterVsIaipSrWL}{35/4}
\newcommand{\SOneDeltaOutageVsIaipSrCI}{[-0.0117, -0.0035]}
\newcommand{\SOneDeltaOutageVsIaipSrRel}{24.2}
\newcommand{\SOneDeltaOutageVsIaipSrSign}{-}
\newcommand{\SOneDeltaOutageVsIaipSrWL}{25/5}
\newcommand{\SOneGapCandidate}{.0159}
\newcommand{\SOneGapIaipSr}{.0231}
\newcommand{\SOneGapPortfolioHinted}{.0287}
\newcommand{\SOneGapRandomRounds}{.0206}
\newcommand{\SOneLocalS}{1.03}
\newcommand{\SOneMeanCandidate}{.0255}
\newcommand{\SOneMeanIaipSr}{.0340}
\newcommand{\SOneMeanPortfolioHinted}{.0446}
\newcommand{\SOneMeanRandomRounds}{.0310}
\newcommand{\SOneNonInferior}{yes}
\newcommand{\SOneNotValidated}{0}
\newcommand{\SOneRandomRoundsVsIaipSrCI}{[-0.0045, -0.0015]}
\newcommand{\SOneRandomRoundsVsIaipSrGapCI}{[-0.0040, -0.0012]}
\newcommand{\SOneRandomRoundsVsIaipSrGapRel}{10.9}
\newcommand{\SOneRandomRoundsVsIaipSrGapSign}{-}
\newcommand{\SOneRandomRoundsVsIaipSrGapWL}{31/9}
\newcommand{\SOneRandomRoundsVsIaipSrHolm}{8.4\cdot10^{-6}}
\newcommand{\SOneRandomRoundsVsIaipSrRel}{8.6}
\newcommand{\SOneRandomRoundsVsIaipSrSign}{-}
\newcommand{\SOneRandomRoundsVsIaipSrWL}{34/6}
\newcommand{\SOneRecords}{8{,}800}
\newcommand{\SOneRoundsImproving}{2.0}
\newcommand{\SOneRoundsMean}{2.9}
\newcommand{\SOneSeedsMedian}{35}
\newcommand{\SOneSizeOneFiveVsIaipSrCI}{[-0.0060, -0.0036]}
\newcommand{\SOneSizeOneFiveVsIaipSrRel}{17.8}
\newcommand{\SOneSizeOneFiveVsIaipSrSign}{-}
\newcommand{\SOneSizeOneFiveVsIaipSrWL}{20/0}
\newcommand{\SOneSizeTwoZeroVsIaipSrCI}{[-0.0159, -0.0083]}
\newcommand{\SOneSizeTwoZeroVsIaipSrRel}{29.4}
\newcommand{\SOneSizeTwoZeroVsIaipSrSign}{-}
\newcommand{\SOneSizeTwoZeroVsIaipSrWL}{19/1}
\newcommand{\SOneUnits}{40}
\newcommand{\SSevenCandidateVsIaipSrCI}{[-0.0094, -0.0050]}
\newcommand{\SSevenCandidateVsIaipSrGapCI}{[-0.0066, -0.0029]}
\newcommand{\SSevenCandidateVsIaipSrGapRel}{25.0}
\newcommand{\SSevenCandidateVsIaipSrGapSign}{-}
\newcommand{\SSevenCandidateVsIaipSrGapWL}{27/7}
\newcommand{\SSevenCandidateVsIaipSrHolm}{7.7\cdot10^{-7}}
\newcommand{\SSevenCandidateVsIaipSrRel}{22.4}
\newcommand{\SSevenCandidateVsIaipSrSign}{-}
\newcommand{\SSevenCandidateVsIaipSrWL}{31/3}
\newcommand{\SSevenCandidateVsPortfolioHintedCI}{[-0.0228, -0.0147]}
\newcommand{\SSevenCandidateVsPortfolioHintedGapCI}{[-0.0156, -0.0086]}
\newcommand{\SSevenCandidateVsPortfolioHintedGapRel}{46.0}
\newcommand{\SSevenCandidateVsPortfolioHintedGapSign}{-}
\newcommand{\SSevenCandidateVsPortfolioHintedGapWL}{33/1}
\newcommand{\SSevenCandidateVsPortfolioHintedHolm}{2.3\cdot10^{-10}}
\newcommand{\SSevenCandidateVsPortfolioHintedRel}{43.0}
\newcommand{\SSevenCandidateVsPortfolioHintedSign}{-}
\newcommand{\SSevenCandidateVsPortfolioHintedWL}{34/0}
\newcommand{\SSevenGapCandidate}{.0141}
\newcommand{\SSevenGapIaipSr}{.0188}
\newcommand{\SSevenGapPortfolioHinted}{.0260}
\newcommand{\SSevenMeanCandidate}{.0248}
\newcommand{\SSevenMeanIaipSr}{.0319}
\newcommand{\SSevenMeanPortfolioHinted}{.0434}
\newcommand{\SSevenNonInferior}{yes}
\newcommand{\SSevenNotValidated}{0}
\newcommand{\SSevenRecords}{5{,}610}
\newcommand{\SSevenUnits}{34}
\newcommand{\SSixCandidateVsIaipSrCI}{[-0.0060, -0.0031]}
\newcommand{\SSixCandidateVsIaipSrGapCI}{[-0.0036, -0.0010]}
\newcommand{\SSixCandidateVsIaipSrGapRel}{18.9}
\newcommand{\SSixCandidateVsIaipSrGapSign}{-}
\newcommand{\SSixCandidateVsIaipSrGapWL}{17/3}
\newcommand{\SSixCandidateVsIaipSrHolm}{4.0\cdot10^{-5}}
\newcommand{\SSixCandidateVsIaipSrRel}{21.5}
\newcommand{\SSixCandidateVsIaipSrSign}{-}
\newcommand{\SSixCandidateVsIaipSrWL}{19/1}
\newcommand{\SSixCandidateVsPortfolioHintedCI}{[-0.0181, -0.0103]}
\newcommand{\SSixCandidateVsPortfolioHintedGapCI}{[-0.0092, -0.0035]}
\newcommand{\SSixCandidateVsPortfolioHintedGapRel}{38.5}
\newcommand{\SSixCandidateVsPortfolioHintedGapSign}{-}
\newcommand{\SSixCandidateVsPortfolioHintedGapWL}{17/3}
\newcommand{\SSixCandidateVsPortfolioHintedHolm}{3.8\cdot10^{-6}}
\newcommand{\SSixCandidateVsPortfolioHintedRel}{46.0}
\newcommand{\SSixCandidateVsPortfolioHintedSign}{-}
\newcommand{\SSixCandidateVsPortfolioHintedWL}{20/0}
\newcommand{\SSixGapCandidate}{.0098}
\newcommand{\SSixGapIaipSr}{.0121}
\newcommand{\SSixGapPortfolioHinted}{.0159}
\newcommand{\SSixMeanCandidate}{.0164}
\newcommand{\SSixMeanIaipSr}{.0209}
\newcommand{\SSixMeanPortfolioHinted}{.0305}
\newcommand{\SSixNonInferior}{yes}
\newcommand{\SSixNotValidated}{0}
\newcommand{\SSixRecords}{1{,}980}
\newcommand{\SSixUnits}{20}
\newcommand{\SThreeCandidateVsIaipSrCI}{[-0.0026, -0.0002]}
\newcommand{\SThreeCandidateVsIaipSrGapCI}{[-0.0028, -0.0006]}
\newcommand{\SThreeCandidateVsIaipSrGapRel}{18.7}
\newcommand{\SThreeCandidateVsIaipSrGapSign}{-}
\newcommand{\SThreeCandidateVsIaipSrGapWL}{13/7}
\newcommand{\SThreeCandidateVsIaipSrHolm}{0.12}
\newcommand{\SThreeCandidateVsIaipSrRel}{7.9}
\newcommand{\SThreeCandidateVsIaipSrSign}{-}
\newcommand{\SThreeCandidateVsIaipSrWL}{14/6}
\newcommand{\SThreeCandidateVsPortfolioHintedCI}{[-0.0045, -0.0014]}
\newcommand{\SThreeCandidateVsPortfolioHintedGapCI}{[-0.0046, -0.0015]}
\newcommand{\SThreeCandidateVsPortfolioHintedGapRel}{29.6}
\newcommand{\SThreeCandidateVsPortfolioHintedGapSign}{-}
\newcommand{\SThreeCandidateVsPortfolioHintedGapWL}{15/5}
\newcommand{\SThreeCandidateVsPortfolioHintedHolm}{0.024}
\newcommand{\SThreeCandidateVsPortfolioHintedRel}{15.6}
\newcommand{\SThreeCandidateVsPortfolioHintedSign}{-}
\newcommand{\SThreeCandidateVsPortfolioHintedWL}{16/4}
\newcommand{\SThreeGapCandidate}{.0071}
\newcommand{\SThreeGapIaipSr}{.0087}
\newcommand{\SThreeGapPortfolioHinted}{.0101}
\newcommand{\SThreeMeanCandidate}{.0154}
\newcommand{\SThreeMeanIaipSr}{.0167}
\newcommand{\SThreeMeanPortfolioHinted}{.0182}
\newcommand{\SThreeNonInferior}{yes}
\newcommand{\SThreeNotValidated}{0}
\newcommand{\SThreeRecords}{3{,}300}
\newcommand{\SThreeUnits}{20}
\newcommand{\STwoCandidateVsIaipSrCI}{[-0.0081, -0.0029]}
\newcommand{\STwoCandidateVsIaipSrGapCI}{[-0.0069, -0.0021]}
\newcommand{\STwoCandidateVsIaipSrGapRel}{18.7}
\newcommand{\STwoCandidateVsIaipSrGapSign}{-}
\newcommand{\STwoCandidateVsIaipSrGapWL}{26/14}
\newcommand{\STwoCandidateVsIaipSrHolm}{0.038}
\newcommand{\STwoCandidateVsIaipSrRel}{13.2}
\newcommand{\STwoCandidateVsIaipSrSign}{-}
\newcommand{\STwoCandidateVsIaipSrWL}{27/13}
\newcommand{\STwoCandidateVsPortfolioHintedCI}{[-0.0158, -0.0100]}
\newcommand{\STwoCandidateVsPortfolioHintedGapCI}{[-0.0122, -0.0070]}
\newcommand{\STwoCandidateVsPortfolioHintedGapRel}{33.2}
\newcommand{\STwoCandidateVsPortfolioHintedGapSign}{-}
\newcommand{\STwoCandidateVsPortfolioHintedGapWL}{34/6}
\newcommand{\STwoCandidateVsPortfolioHintedHolm}{3.7\cdot10^{-7}}
\newcommand{\STwoCandidateVsPortfolioHintedRel}{26.8}
\newcommand{\STwoCandidateVsPortfolioHintedSign}{-}
\newcommand{\STwoCandidateVsPortfolioHintedWL}{36/4}
\newcommand{\STwoGapCandidate}{.0193}
\newcommand{\STwoGapIaipSr}{.0238}
\newcommand{\STwoGapPortfolioHinted}{.0289}
\newcommand{\STwoMeanCandidate}{.0351}
\newcommand{\STwoMeanIaipSr}{.0405}
\newcommand{\STwoMeanPortfolioHinted}{.0480}
\newcommand{\STwoNonInferior}{yes}
\newcommand{\STwoNotValidated}{0}
\newcommand{\STwoRecords}{6{,}600}
\newcommand{\STwoUnits}{40}


\newtheorem{proposition}{Proposition}
\newtheorem{corollary}{Corollary}
\newcommand{\base}{\textsc{SR}}
\newcommand{\crm}{\textsc{SR+CR}}
\newcommand{\rr}{\textsc{SR+RR}}
\newcommand{\ph}{\textsc{Port+Hint}}

\setcounter{secnumdepth}{0}

\title{Sound Neighbourhood Certificates and Critical-Region Rounds\\for Anytime Rescheduling under a Deadline}
\author{
    Anonymous Submission
}
\affiliations{
}

\begin{document}

\maketitle

\begin{abstract}
After a disruption, a job shop must be rescheduled within a wall-clock budget, and a common pipeline keeps a validated incumbent, re-optimizes a restricted neighbourhood of the operations the event affected with CP-SAT, and then continues with the unrestricted solver. We ask whether sound information about the improvement still possible inside a neighbourhood can decide when to run the restricted stage, which machine-order decisions to release, and how to split the budget. We give a cheap lower bound for restricted neighbourhoods that handles outages, release dates and commitments, a minimum-cut procedure that releases a minimum set of machine arcs meeting every obstructing chain, and certificate-aware budget rules, and check them against exact oracles on \OracleEnumCases{} enumerated and \OracleCpsatCases{} CP-SAT-solved cases with no violation. On development streams the certificate is rarely decisive, and CP-SAT disposes of non-improving neighbourhoods in milliseconds; skipping, releasing and certificate-aware allocation have no measurable effect, and a budget rule built from non-significant evidence is harmful. What helps is spending the restricted window on repeated rounds seeded at the incumbent's critical region. Frozen on development data and evaluated on 40 held-out streams with five solver seeds, this reduces the primal integral of the pipeline by \SOneCandidateVsIaipSrRel\% (\SOneCandidateVsIaipSrWL{} streams) and beats the same rounds with random seeds by \SOneCandidateVsRandomRoundsRel\%; with 2\,s budgets and with four solver workers its mean primal integral is also lower, with bootstrap intervals that exclude zero, and secondary settings (arrival-only and outage-only streams, 30\,s budgets, streams built on 34 published instances) show the same direction. Critical-region neighbourhoods are a known idea; our contributions are the sound certificate framework, the evidence on when certificate information is and is not decision-relevant under a deadline, and a pre-registered evaluation.
\end{abstract}

\section{Introduction}

Schedules are revised repeatedly as jobs arrive, durations change and machines fail, each time within a time limit \citep{vieira2003,ouelhadj2009}. With a constraint solver, a natural pipeline keeps a validated incumbent built from the previous schedule or from dispatching rules, re-optimizes only a neighbourhood of the operations a disruption affected, and spends the rest of the budget on the unrestricted problem. Such selective repair is used with CP-SAT by recent work that learns which operations to release \citep{linyang2026}, and the choice of rescheduling scope has been studied directly \citep{sun2026}.

Two decisions in such a pipeline look like they should benefit from mathematics rather than heuristics: whether a restricted neighbourhood can improve the incumbent at all, and how much of the deadline the restricted stage deserves. A valid lower bound $L$ on the best schedule inside a neighbourhood bounds the possible improvement by $U-L$, where $U$ is the incumbent's makespan; if $L \ge U$, the neighbourhood is useless. We call such information an \emph{improvement-opportunity certificate} and study whether it changes decisions for the better once its cost is counted.

Our findings are mixed, and we report them as such.
\begin{itemize}
\item The certificate is sound and cheap (Propositions 1--3), verified against exact oracles with no false certificate, and it exposes modeling pitfalls that also affect simpler restricted models (zero-length operations, solver-versus-validator semantics, horizons).
\item As a decision device it is nearly useless here: it is decisive for \PilotCertDecisive{} of \PilotCertSteps{} disruption-centred neighbourhoods, CP-SAT proves the same non-improvement within a median of \DiagInfeasibleMs\,ms, and certificate-aware skipping, release and allocation have no measurable effect on the primal integral (Table~\ref{tab:dev}).
\item A deadline rule derived from earlier non-significant results (skip the restricted stage with short budgets or several workers) is harmful, a reminder that absence of significance is not evidence of absence.
\item What does help is to keep the restricted window but fill it with repeated restricted rounds, each seeded at the operations of the current incumbent's longest binding chains and expanded by breadth-first search. Under a configuration frozen on development streams, this improves the pipeline on held-out streams in all three predeclared settings (Table~\ref{tab:conf}), and a control with random seeds shows that \emph{where} the rounds search matters.
\end{itemize}
Critical-path neighbourhoods are classical in job-shop local search \citep{vanlaarhoven1992,nowicki1996,balas1998}, so we do not claim the rounds as a new idea; we claim the evidence.

\section{Related Work}

\paragraph{Selective rescheduling.} Affected-operations rescheduling re-optimizes only the operations a disruption affects \citep{abumaizar1997,subramaniam2003}. \citet{sun2026} compare optimization scopes from minimal repair to global re-optimization in dynamic flexible job shops. SLR2 \citep{linyang2026} encodes the incumbent and the disturbance in a heterogeneous graph, uses a graph neural network to split operations into kept and released sets, and repairs the released neighbourhood with a reduced CP-SAT model under a time budget. Learning which variables to fix is also the basis of rolling-horizon methods for flexible job shops \citep{li2025lrho,li2026graphrho}, the latter using critical-path information as a training bias. Our release rule is not learned: it is a structural necessary condition computed per instance, and we evaluate anytime quality under one end-to-end deadline.

\paragraph{Critical paths, bounds and neighbourhoods.} In the disjunctive graph of a complete selection the makespan equals the longest path \citep{balas1969}; neighbourhoods that reverse critical arcs or move operations within critical blocks underlie job-shop local search \citep{vanlaarhoven1992,nowicki1996,balas1998}, and one-machine relaxations give classical lower bounds \citep{carlier1982}. Our certificate is a relaxation bound of this kind for a \emph{restricted} problem, and our release step uses the max-flow min-cut theorem \citep{fordfulkerson1956,edmondskarp1972}.

\paragraph{CP-based LNS and CP-SAT.} Large neighbourhood search re-solves part of an incumbent with constraint programming \citep{shaw1998,godard2005,pacino2011} and descends from the shuffle technique for job shops \citep{applegate1991}. CP-SAT \citep{cpsat,perron2023} runs its own scheduling neighbourhoods when given several workers; with one worker it runs a single search unless its experimental interleaved mode is enabled. Anytime quality is measured by the primal integral \citep{berthold2013}.

\section{Problem and Base Pipeline}

A \emph{stream} is a sequence of job-shop instances $I_0,\dots,I_T$; $I_t$ arises from $I_{t-1}$ by a job arrival, a cancellation, a duration change or a machine outage. Each instance has integer durations $p_v \ge 0$, release dates on first operations, disjoint outage intervals per machine, and possibly committed start times. Each instance is solved within a wall-clock budget $B$ covering construction, model building, solving and validation, and a method receives the schedule it returned for $I_{t-1}$. With $z(\tau)$ the best makespan found by time $\tau$ and $z^*$ the best makespan any compared run found for the instance, the gap is $g(\tau)=(z(\tau)-z^*)/z^*$, $g=1$ before the first schedule, and the primal integral is $\mathrm{PI}=\frac1B\int_0^B g(\tau)\,d\tau$.

The base pipeline \base{} (i) builds a portfolio of validated incumbents---repairs of the previous schedule and dispatching rules---and retains each improving one when it is built; (ii) restricts CP-SAT to a neighbourhood $N$ grown by breadth-first search from the operations whose data changed, capped at 25\% of all operations: every pair of operations outside $N$ that is consecutive in the incumbent's machine order keeps its order, and $C_{\max} \le U-1$ is imposed; this stage ends at elapsed time $0.6B$ or when CP-SAT proves the restricted problem optimal or infeasible; (iii) runs unrestricted CP-SAT for the rest of the budget, hinted with the current incumbent. \ph{} is \base{} without stage (ii).

\section{Improvement-Opportunity Certificates}

Let $S$ be the incumbent, $U=C_{\max}(S)$, $N$ a set of free operations and $R$ a set of released arcs. The restricted problem $P(S,N,R)$ adds to the instance the precedence $s_v \ge s_u+p_u$ for every \emph{kept arc}: a pair consecutive in $S$'s machine order among the operations of \emph{positive} duration outside $N$, unless it is in $R$. $C^*(N,R)$ is its optimum.

Let $G$ be the graph of job arcs and kept arcs. In a topological order, let $\mathrm{ready}(v)$ be the maximum of $v$'s release date (if it is a first operation) and $h(u)+p_u$ over predecessors $u$; let $h(v)$ be the commitment $f_v$ if $v$ is committed, the earliest outage-free start at or after $\mathrm{ready}(v)$ if $p_v>0$, and $\mathrm{ready}(v)$ otherwise. Define $L_{\mathrm{path}}=\max_v h(v)+p_v$, $L_{\mathrm{mach}}=\max_m \bigl(\min_{v\in m} h(v)+\sum_{v\in m} p_v+\min_{v\in m} q(v)\bigr)$ with $q$ the longest calendar-free path after $v$ in $G$, and $L=\max(L_{\mathrm{path}},L_{\mathrm{mach}})$.

\begin{proposition}
Every schedule of $P(S,N,R)$ has $s_v\ge h(v)$ for all $v$ and makespan at least $L$; hence $L \le C^*(N,R)$.
\end{proposition}
\begin{proof}
By induction over the topological order: predecessors start no earlier than their bounds, so $s_v \ge \mathrm{ready}(v)$; a committed $v$ starts at $f_v$; for $p_v>0$, $s_v$ is an outage-free start at or after $\mathrm{ready}(v)$, and the earliest such start is a monotone function of $\mathrm{ready}(v)$. For $L_{\mathrm{mach}}$, the operations of a machine run without overlap from at least $\min h$, and the last one still needs its tail.
\end{proof}

\begin{corollary}
$S\in P(S,N,R)$, so $C^*(N,R)\le U$; if $L\ge U$, no schedule in the neighbourhood is better than $S$. With $N$ the set of all operations, $P$ is the full problem and $L\ge U$ proves $S$ optimal.
\end{corollary}
Membership needs the positive-duration restriction: a zero-length operation at the start of another operation's interval would receive an arc $S$ violates. These statements are elementary relaxation arguments; we do not present them as new results.

\paragraph{Releasing obstructions.} If $L_{\mathrm{path}}=U$, call an arc $(u,v)$ \emph{binding} if $h(u)+p_u=\mathrm{ready}(v)=h(v)$, and call $v$ a \emph{chain start} if its bound would not drop if all its incoming arcs were removed (committed, held by an outage, or attained by its release date or 0, ties included). Every path of length $U$ contains a binding chain from a chain start to an operation $t$ with $h(t)+p_t=U$.
\begin{proposition}
In the network of binding arcs, with capacity $w(e)$ for releasable machine arcs, $\infty$ for job arcs, a source joined to all chain starts and all such $t$ joined to a sink, a minimum cut is a minimum-weight set of releasable arcs that meets every obstructing chain; an infinite cut means some chain contains no releasable arc.
\end{proposition}
\begin{proposition}
If, after releasing a finite minimum cut, no operation's bound becomes held by an outage at an unchanged value, then $L_{\mathrm{path}}<U$ for the new neighbourhood.
\end{proposition}
The condition can fail: after a release, an operation may be held by an outage exactly at its old bound, so the procedure recomputes the certificate and repeats, and only the recomputed certificate is ever used. A non-decisive certificate does not mean that an improving schedule exists.

\paragraph{Pitfalls found by the oracles.} Three modeling details matter for any restricted model: chaining zero-length operations by start order can exclude the incumbent; CP-SAT's no-overlap rejects a zero-length interval strictly inside an outage although a validator that ignores zero-length intervals accepts it; and an end-time horizon of total duration plus last outage end can exclude feasible schedules once release dates are present. None occurs in our experimental instances.

\paragraph{Cost.} Computing $L$ and the cut is linear in the graph size apart from the max-flow; on our instances the certificate takes \PilotCertMs\,ms and the release \PilotReleaseMs\,ms on average.

\section{Certificate-Guided Decisions and Critical-Region Rounds}

We tested four uses of the certificate inside \base: \emph{skip} a neighbourhood certified non-improving; \emph{release} obstructions of a certified neighbourhood (Proposition 2); \emph{allocate} a short restricted stage ($0.2B$) when the certified opportunity $(U-L)/U$ is below 2\%; and a \emph{budget rule} that skips the restricted stage with fewer than 4\,s or at least four workers (derived from earlier exploratory results).

\paragraph{Rounds.} After the restricted stage ends early---as it usually does---the critical-region variant \crm{} runs further restricted rounds inside the same window, while the previous round improved and at most ten in total. Each round recomputes the binding chains of the current incumbent's own machine orders (of its left-shifted version if it has idle time), takes every operation on them as a seed, grows a neighbourhood by the same breadth-first expansion (25\% of operations) and solves it with $C_{\max}\le U-1$. Two controls differ only in the seeds: the endpoints of a minimum cut of those chains, and a random set of operations of the same size as the cut seeds (\rr). All stages share one monotonic deadline; the unrestricted stage is a separate CP-SAT run that receives only the incumbent as a hint.

\begin{table*}[t]
\centering
\small
\setlength{\tabcolsep}{5pt}
\begin{tabular}{lrrrrrr}
\toprule
Setting & $n$ & \crm & \base & \ph & vs \base & vs \ph \\
\midrule
\csname @@input\endcsname tables/conf.tex 
\bottomrule
\end{tabular}
\caption{Held-out confirmation (stream seeds 901--920, solver seeds 0--4; configuration frozen on development data). Mean PI per method; relative difference of \crm{} and streams won/lost against each comparator. $n$: streams.}
\label{tab:conf}
\end{table*}

\section{Experimental Methodology}

Instances are full $m{\times}m$ job shops ($m\in\{15,20\}$) with durations uniform in $[5,50]$; ten deltas per stream are drawn with weights 0.4 (arrival), 0.2 (cancellation), 0.25 (duration change) and 0.15 (outage). CP-SAT (OR-Tools 9.15) runs with one worker and $B=8$\,s unless stated; each process is pinned to its own physical core of an Intel Core i7-12700 (four physical cores per process with four workers); methods of a stream run sequentially in one process in shuffled order. Every returned schedule is checked by an independent validator; in an audit of \AuditSolutions{} intermediate CP-SAT solutions, \AuditInvalid{} were invalid and all had the reported makespan.

\paragraph{Development and freeze.} All method design used development streams (stream seeds 301--310) in three pilot rounds with solver seeds 0, 1 and 2, each round's question and decision rule written down before it ran. The configuration was then frozen with its comparators, settings and success criteria in a file whose hash the analysis script checks.

\paragraph{Confirmation.} The held-out data are stream seeds 901--920 (40 streams, never used before) with solver seeds 0--4. Settings: S1, $B=8$\,s, one worker, with \crm, \base, \ph{} and \rr; S2, $B=2$\,s; S3, four workers, 20 streams; and, as secondary settings run afterwards in a fixed order, arrival-only and outage-only streams (20 each), $B=30$\,s (20 streams, solver seeds 0--2), and streams built on the published instances ta01--ta20 \citep{taillard1993}, dmu01--05 \citep{demirkol1998}, swv01--05 \citep{storer1992} and yn1--4 \citep{yamada1992}, one stream each, with the base instance as the unit. The unit of inference is the stream: PI is averaged over steps 1--10 and seeds within a stream, $z^*$ is shared per physical instance, and comparisons use the exact sign test over streams with Holm correction and a stream-clustered bootstrap 95\% interval (10{,}000 resamples). Criteria fixed in advance: \crm{} beats \base{} by at least 5\% in S1 with an interval excluding zero; beats \ph{} by at least 10\%; beats \rr; is non-inferior to \base{} in S2 and S3 with a margin of 5\% of \base's mean PI; shows no final-gap regression; and every output is validated.

\section{Results}

\paragraph{Soundness.} Against the enumeration oracle (every machine order consistent with the kept arcs, scheduled by an independent earliest-start routine) the bound was valid in all \OracleEnumCases{} cases, of which \OracleEnumDecisive{} were decisive, each confirmed; against CP-SAT optima of the restricted model, including zero-length operations, in all \OracleCpsatCases{} cases (\OracleCpsatDecisive{} decisive); the global certificate in all \OracleGlobalCases{} cases. On \OracleCutCases{} decisive cases the minimum cut met every obstructing chain and was a minimum hitting set by enumeration (\OracleCutMisses{} exceptions).

\paragraph{What the certificate changes.} Replaying the development runs of \base{} offline, the certificate was decisive for \DiagDecisive{} of \DiagSteps{} restricted stages and explained \DiagExplained{} of the \DiagInfeasible{} stages that CP-SAT proved non-improving; those proofs took a median of \DiagInfeasibleMs\,ms, and stages that found the restricted optimum a median of \DiagOptimalS\,s, far below the $0.6B$ window. Skipping or shortening the stage can therefore save only milliseconds, and Table~\ref{tab:dev} confirms that the certificate-based variants are indistinguishable from \base{}, although the allocation rule changed the decision in \PilotShortStage{} of \PilotCertSteps{} steps. The budget rule, in contrast, was harmful: \base's restricted stage helps with 2\,s budgets and with four workers, and skipping it lost $\PilotTwoAFourVsAOneRel$\% and $\PilotFourAFourVsAOneRel$\%. The rule had been derived from comparisons that were merely not significant.

\begin{table*}[t]
\centering
\small
\setlength{\tabcolsep}{5pt}
\begin{tabular}{llrrr}
\toprule
Variant & Pair & $\Delta$PI & W/L & 95\% CI \\
\midrule
\csname @@input\endcsname tables/dev.tex 
\bottomrule
\end{tabular}
\caption{Development pilots (exploratory; 20 streams at 8\,s, 10 at 2\,s and with four workers; one solver seed each). A1 = \base, A0 = \ph, A4 = budget rule, A9 = \ph{} with CP-SAT's interleaved search; C1, C2, C3 = rounds seeded at min-cut endpoints, at the critical region, at random operations. $\Delta$PI: relative difference of mean PI; W/L: streams won/lost by the first method; CI of the mean stream difference.}
\label{tab:dev}
\end{table*}

\paragraph{Rounds.} Rounds seeded at the min-cut endpoints were not distinguishable from rounds with random seeds, so the min-cut seeding is not credited; rounds seeded at the whole critical region improved \base{} on two development seeds (Table~\ref{tab:dev}), and that variant was frozen.

\paragraph{Confirmation.} Table~\ref{tab:conf} reports the held-out settings. In S1, \crm{} reduces mean PI by \SOneCandidateVsIaipSrRel\% relative to \base{} (\SOneCandidateVsIaipSrWL{} streams, interval $\SOneCandidateVsIaipSrCI$, Holm $p=\SOneCandidateVsIaipSrHolm$), by \SOneCandidateVsPortfolioHintedRel\% relative to \ph{} (\SOneCandidateVsPortfolioHintedWL), and by \SOneCandidateVsRandomRoundsRel\% relative to \rr{} (\SOneCandidateVsRandomRoundsWL); random-seed rounds themselves improve \base{} by \SOneRandomRoundsVsIaipSrRel\% (\SOneRandomRoundsVsIaipSrWL). The final gap at $B$ improves by \SOneCandidateVsIaipSrGapRel\% (\SOneCandidateVsIaipSrGapWL). With 2\,s budgets \crm{} is \STwoCandidateVsIaipSrRel\% better than \base{} (\STwoCandidateVsIaipSrWL, interval $\STwoCandidateVsIaipSrCI$) and with four workers \SThreeCandidateVsIaipSrRel\% better (\SThreeCandidateVsIaipSrWL, interval $\SThreeCandidateVsIaipSrCI$, Holm $p=\SThreeCandidateVsIaipSrHolm$), so the predeclared non-inferiority criteria hold; in the four-worker setting the sign test is not significant after Holm correction, although the bootstrap interval excludes zero. The secondary settings agree: arrival-only $\SFourCandidateVsIaipSrSign\SFourCandidateVsIaipSrRel$\% (\SFourCandidateVsIaipSrWL), outage-only $\SFiveCandidateVsIaipSrSign\SFiveCandidateVsIaipSrRel$\% (\SFiveCandidateVsIaipSrWL), 30\,s $\SSixCandidateVsIaipSrSign\SSixCandidateVsIaipSrRel$\% (\SSixCandidateVsIaipSrWL), and published bases $\SSevenCandidateVsIaipSrSign\SSevenCandidateVsIaipSrRel$\% (\SSevenCandidateVsIaipSrWL{} base instances). In S1 the reduction is similar across delta types (arrivals $\SOneDeltaArrivalVsIaipSrSign\SOneDeltaArrivalVsIaipSrRel$\%, cancellations $\SOneDeltaCancellationVsIaipSrSign\SOneDeltaCancellationVsIaipSrRel$\%, duration changes $\SOneDeltaDurationJitterVsIaipSrSign\SOneDeltaDurationJitterVsIaipSrRel$\%, outages $\SOneDeltaOutageVsIaipSrSign\SOneDeltaOutageVsIaipSrRel$\%; exploratory) and larger for $20{\times}20$ ($\SOneSizeTwoZeroVsIaipSrSign\SOneSizeTwoZeroVsIaipSrRel$\%) than for $15{\times}15$ ($\SOneSizeOneFiveVsIaipSrSign\SOneSizeOneFiveVsIaipSrRel$\%). All \SOneRecords{} S1 runs and every run of the other settings returned a validated schedule. On average \crm{} ran \SOneRoundsMean{} extra rounds per instance, \SOneRoundsImproving{} of them improving, seeded at a median of \SOneSeedsMedian{} operations, and used \SOneLocalS\,s of the 4.8\,s window, leaving the rest to the unrestricted solver.



\section{Limitations}

The instances are synthetic full job shops of two sizes with one generator; development and held-out streams come from the same generator. $z^*$ is the best value found, not a proven optimum. Only makespan is optimized and stability is not controlled. The development pilots used one solver seed per round and showed noticeable seed dependence: the first rounds variant improved \base{} by \PilotOneASixrVsAOneRel\% on solver seed 0 but by \PilotEightASixrVsAOneRel\% on seed 1. The four-worker setting used 20 streams, and its sign test is not significant after correction. CP-SAT's interleaved search, which matched \ph{} at 8\,s in development, was not compared with \crm{} on held-out data; whether it closes the gap with longer budgets is open. CP-SAT's internal search changes with its version. The certificate's negative result is specific to a strong complete solver: with a weaker restricted-stage solver, a millisecond proof of non-improvement would not be free, and the certificate might matter.

\section{Conclusion}

Sound neighbourhood certificates are easy to get right only with care---the oracles found three modeling pitfalls---and, once right, they carry information CP-SAT already extracts in milliseconds, so they did not improve any decision we tested. Under the same deadline, the time is better spent on repeated restricted rounds around the incumbent's critical region, which improved a strong selective-reoptimization pipeline on held-out streams in every predeclared setting.

\bibliography{references}

\end{document}
